% LAD-DA3: Height-Anchored Depth for Roadside V2X Scene Reconstruction
%
% Drafted in the CVPR-paper style. The official cvpr.sty is not vendored;
% drop it next to this file before building. Falls back to a plain article
% class so the source is still readable / compilable for quick iteration.

\IfFileExists{cvpr.sty}{
    \documentclass[10pt,twocolumn,letterpaper]{article}
    \usepackage{cvpr}
}{
    % Fallback for quick local iteration without the CVPR kit.
    \documentclass[10pt,twocolumn,letterpaper]{article}
    \usepackage[margin=0.75in]{geometry}
    \usepackage{times}
    \newcommand{\cvprfinalcopy}{}
    \newcommand{\confName}{CVPR}
    \newcommand{\confYear}{20XX}
}

\usepackage{amsmath, amssymb, amsfonts}
\usepackage{graphicx}
\usepackage{booktabs}
\usepackage{multirow}
\usepackage{subcaption}
\usepackage{xcolor}
\usepackage{algorithm}
\usepackage[noend]{algpseudocode}
\usepackage[hidelinks]{hyperref}
\usepackage{microtype}

% Compact in-line TODOs, easy to grep.
\newcommand{\todo}[1]{\textcolor{red}{[TODO: #1]}}
% A short name for the system; saves typing and keeps it consistent.
\newcommand{\ourmethod}{\textsc{Lad-DA3}\xspace}
% xspace if not provided by cvpr.sty:
\IfFileExists{xspace.sty}{\usepackage{xspace}}{\newcommand{\xspace}{}}

% \cvprfinalcopy   % uncomment for the camera-ready
\def\paperID{XXXX}
\def\confName{CVPR}
\def\confYear{20XX}

\title{Height-Anchored Depth: Closed-Form Foundation Depth Calibration\\
for Roadside V2X Scene Reconstruction}

\author{Anonymous Authors\\
Anonymous Institution\\
{\tt\small anonymous@example.org}
}

\begin{document}
\maketitle


%======================================================================
\begin{abstract}
Foundation monocular depth models such as Depth Anything 3 produce
relative depth maps that are remarkably faithful in shape but suffer
from \emph{scale drift} in out-of-distribution settings — most notably
the roadside vehicle-to-everything (V2X) configuration, where cameras
are mounted 5–10\,m above the road plane and observe traffic 50–200\,m
away. The standard fix is to align the predicted depth to LiDAR returns
via a global scale or a per-pixel affine map; both rely on hundreds of
thousands of noisy per-pixel anchors and both collapse at long range,
exactly where roadside applications need the answer.

We argue that the right primitive for roadside scenes is the
\emph{world-frame height} of segmented instances, not per-pixel camera
depth. From a Vehicle-to-Everything (V2X) bounding box we read off two
height anchors — the road surface and the object top — and the
predicted relative depth becomes metric via a one-line linear system,
solved in closed form per camera and per instance. We call this formulation
\textbf{Height-Anchored Depth} (HAD), and its axis-aligned variant
\textbf{AA-HAD}. Around this core we build an end-to-end pipeline that
combines AA-HAD with a LiDAR-derived 2.5D ground grid, a LiDAR-priority
voxel fusion, a per-object accumulation in object-local frames, and a
trained DGCNN+conditional-flow-matching residual head that refines the
closed-form prediction. The full system is shipped as a typed CLI
(\texttt{lad}) with multi-GPU sharding and an
\texttt{outputs/<scene>/<stage>/<run>/} discovery convention that turns
the whole eight-stage pipeline into one command.

We demonstrate \ourmethod{} on the TianJin Roadside V2X dataset. A
model trained on one intersection transfers zero-shot to an unseen
intersection without retraining, producing dense static reconstructions
that are visibly cleaner than the closed-form baseline and dynamic-object
trajectories that move smoothly across bird's-eye-view video frames.
\end{abstract}


%======================================================================
\section{Introduction}
\label{sec:intro}

Foundation monocular depth models~\cite{depthanything3,zoedepth,midas3}
have closed the gap between research-grade and deployment-grade depth
estimation on a wide spectrum of in-the-wild imagery — driving scenes,
indoor scenes, hand-held captures. Two limitations remain stubborn,
however, when the camera moves out of the training distribution. First,
the predictions are typically \emph{relative} (or, in the metric heads,
biased on scenes the model was not trained on), so a downstream
calibration step is needed before the depth can be unprojected into a
metric 3D scene. Second, the standard calibration step — align the
predicted depth to LiDAR returns via a global scale or a per-pixel
affine~\cite{lidaraligndepth} — fails on \emph{long-range} returns, the
exact regime in which roadside Vehicle-to-Everything (V2X)
applications operate.

The roadside V2X setting (Fig.~\ref{fig:teaser}) is special in three
ways relevant to depth estimation. (i) The camera is mounted 5–10\,m
above the road plane and observes traffic at 20–200\,m, so per-pixel
camera-axis depth varies across nearly two orders of magnitude in a
single image. (ii) The available LiDAR is sparse at the same long
ranges where the camera is most uncertain. (iii) V2X annotations
provide \emph{per-instance 3D bounding boxes} for every annotated
vehicle, pedestrian, and roadside fixture — a data product that is
nearly always available in this setting and almost never exploited
by current depth-calibration pipelines.

\paragraph{Key insight.} Per-pixel \emph{camera-axis depth} is the
wrong primitive for a roadside camera. The right primitive is
\emph{world-frame height}: the road surface is approximately level,
each annotated instance has a known vertical extent, and the camera's
pose with respect to the world frame is provided by the calibration
pack. From an instance's 3D bounding box we obtain two height anchors
— the road surface ($Z_{\min}^i$) and the object top ($Z_{\max}^i$).
A simple closed-form linear system maps the foundation model's relative
depth to a metric scale that places those two anchors at the right
heights. With 2 anchors per instance, a single least-squares solve per
camera, and no iteration, the camera depth becomes metric.

\paragraph{Contributions.} We introduce \ourmethod{}, a layered pipeline
that combines this Height-Anchored Depth (HAD) formulation with the
other components needed to produce a dense, multi-camera, multi-frame
3D scene reconstruction in the roadside V2X setting. Concretely:

\begin{itemize}
\setlength{\itemsep}{1pt}
\item We derive the \textbf{HAD calibration} (Sec.~\ref{sec:method-had})
and its axis-aligned approximation \textbf{AA-HAD}
(Sec.~\ref{sec:method-aahad}), giving a one-shot per-camera linear
calibration of a foundation depth model against a handful of V2X
bounding boxes.
\item We design a \textbf{layered system}
(Sec.~\ref{sec:method-pipeline}) — Layer~3 ground-snap on a
LiDAR-derived 2.5D height grid, Layer~4 LiDAR-priority voxel fusion,
and Layer~5 per-object accumulation in object-local frames — that turns
per-frame metric depth into a static-scene reconstruction with
smoothly-placed dynamic actors.
\item We train a \textbf{point-cloud residual flow head}
(Sec.~\ref{sec:method-residual}) — a DGCNN backbone with a conditional
flow-matching head and segment-aware EdgeConv masking — that refines
the closed-form prediction near object edges and on poorly-observed
surfaces.
\item We ship the full system as a \textbf{typed multi-stage CLI}
(\texttt{lad}, Sec.~\ref{sec:impl}) with multi-GPU sharding and a
convention-over-configuration output layout, so a fresh scene goes from
raw frames to a BEV video with a single shell command.
\item We provide a \textbf{dual labels-source} mechanism
(Sec.~\ref{sec:impl-dual}) that uses hand-labeled 1\,Hz V2X bounding
boxes for static-scene accumulation (high spatial precision) and
interpolated 10\,Hz boxes for per-object accumulation and video
rendering (smooth motion). The static-only cloud loses the "trail of
leaked dynamic-object points" artefact that plagues purely-interpolated
pipelines (Sec.~\ref{sec:exp-ablation}).
\end{itemize}

We demonstrate \ourmethod{} qualitatively on the TianJin Roadside V2X
dataset (Sec.~\ref{sec:exp}). A model trained on one intersection
transfers zero-shot to an unseen intersection; the refined point
clouds are visibly cleaner than the closed-form baseline; and the
generated BEV videos show dynamic objects moving smoothly across
frames.

\begin{figure}[t]
  \centering
  \includegraphics[width=\linewidth]{figures/teaser.png}
  \caption{\textbf{Roadside V2X reconstruction from \ourmethod{}.}
  Top-down BEV of a busy intersection. Dense road surface and roadside
  geometry come from a single foundation depth model (Depth Anything
  3) calibrated by per-instance heights from V2X bounding boxes;
  per-object accumulated reconstructions are placed in the scene at
  the current timestamp with linear pose interpolation. No per-pixel
  LiDAR alignment is used. Insets: zoom-ins on a vehicle and a
  traffic-light pole.}
  \label{fig:teaser}
\end{figure}


%======================================================================
\section{Related Work}
\label{sec:related}

\paragraph{Foundation monocular depth.} Depth Anything 3
\cite{depthanything3} and its predecessor Depth Anything~v2
\cite{depthanything2} have set the standard for in-the-wild relative
depth prediction by pre-training on $\sim$$10^7$ images with a
combination of curated synthetic depth and self-distilled
pseudo-labels. ZoeDepth \cite{zoedepth}, Marigold \cite{marigold}, and
the recent generation of diffusion-based depth models target metric
prediction directly, but typically saturate at $\sim$$80$\,m. None of
these models — relative or metric — addresses the roadside V2X regime,
where the depth distribution spans nearly two orders of magnitude in a
single image.

\paragraph{Depth-to-LiDAR alignment.} A long line of work aligns
predicted depth to sparse LiDAR returns via a global scale-and-shift
\cite{midas3}, a per-region affine \cite{regionaffinedepth}, or a
learned per-pixel correction head \cite{learneddepthcalibration}. All
of these treat per-pixel camera-axis depth as the primitive being
aligned. \ourmethod{} differs in two ways: it aligns to per-instance
world-frame heights (not per-pixel camera-axis depth), and the per-
camera solve is closed-form (no iteration, no learned alignment head).

\paragraph{V2X 3D scene reconstruction.} BEVFormer~\cite{bevformer},
BEVFusion~\cite{bevfusion}, and the V2X-Sim/OPV2V family
\cite{v2xsim,opv2v} construct 3D feature volumes from multi-view
roadside imagery, then decode bird's-eye-view object detections.
These pipelines target perception (detection, tracking), not dense
scene reconstruction; their depth representation is implicit and
internal. \ourmethod{} produces an explicit dense world-frame point
cloud that is suitable for both perception and downstream simulation.

\paragraph{Point-cloud refinement / completion.} PCN \cite{pcn},
PoinTr \cite{pointr}, and the PointTransformer~v3
\cite{pointtransformerv3} family refine partial point clouds via a
learned encoder–decoder. \ourmethod{}'s residual head
(Sec.~\ref{sec:method-residual}) uses a DGCNN backbone with a
conditional flow-matching head; the output is a per-point Euclidean
displacement applied to the AA-HAD prediction. This is closer in
spirit to denoising than to completion: the input cloud is already
metric and approximately correct, and the head's job is to remove
residual structure that the closed-form prior does not capture.

\paragraph{Per-object reconstruction.} The closest line of prior work
is multi-frame per-object aggregation in the autonomous-driving stack
\cite{multiframeperobject}, where bounding-box-cropped LiDAR returns
from each frame are placed into an object-local frame and accumulated.
We adopt the same idea (Sec.~\ref{sec:method-perobject}) and add a
foundation-depth-based dense branch alongside the LiDAR branch.


%======================================================================
\section{Method}
\label{sec:method}

We describe \ourmethod{} in two parts. Sec.~\ref{sec:method-had}–
\ref{sec:method-aahad} derive the closed-form HAD calibration that is
the conceptual core of the system. Sec.~\ref{sec:method-pipeline}–
\ref{sec:method-bev} walk through the eight-stage pipeline that wraps
HAD into an end-to-end dense reconstruction.

\subsection{Notation}
\label{sec:method-notation}

Let $\mathbf{K}\in\mathbb{R}^{3\times 3}$ be a pinhole camera's
intrinsic matrix and $\mathbf{T}_{wc}\in\mathrm{SE}(3)$ its world-from-
camera pose. A pixel $(u,v)$ unprojects to a metric world point via
\begin{align}
\mathbf{r}_{cam}(u,v) &= \mathbf{K}^{-1}\,(u, v, 1)^\top
\label{eq:ray} \\
\mathbf{p}_{cam}(u,v) &= z_{cam}(u,v)\,\mathbf{r}_{cam}(u,v)
\label{eq:campt} \\
\mathbf{p}_{world}(u,v) &= \mathbf{T}_{wc}\,\mathbf{p}_{cam}(u,v).
\label{eq:worldpt}
\end{align}
The foundation depth model produces a relative depth map
$\tilde{d}(u,v)$. The task is to recover the metric camera-axis depth
$z_{cam}(u,v)$ from $\tilde{d}(u,v)$.


\subsection{Height-Anchored Depth (HAD)}
\label{sec:method-had}

The standard global-affine calibration \cite{midas3} writes
$z_{cam} = a\,\tilde{d} + b$ and solves for $(a, b)$ by matching $z_{cam}$
to LiDAR returns at every pixel that has one. Two failure modes are
well known: LiDAR returns are sparse at long ranges, so $(a, b)$ is
dominated by close-range pixels; and the per-pixel camera-axis depth
varies across two orders of magnitude in a single roadside image, so
a single $(a, b)$ pair fits the foreground at the expense of the
horizon (or vice versa).

The HAD reformulation drops the per-pixel matching in favour of
\textbf{per-instance world-frame height matching}. Concretely, for each
V2X bounding-box instance $i$ projected into camera $c$, we know:

\begin{itemize}
\item Its top and bottom world-frame heights $Z_{\max}^i$ and
$Z_{\min}^i$ (from the bbox annotation).
\item Two image-space anchor pixels $(u_{\min}^i, v_{\min}^i)$ and
$(u_{\max}^i, v_{\max}^i)$ taken from a segmentation mask
\cite{sam} restricted to the instance.
\item The relative-depth values $\tilde{d}_{\min}^i, \tilde{d}_{\max}^i$
at those anchor pixels.
\end{itemize}

For each anchor we have a per-pixel constraint that the unprojected
\emph{world} height matches the annotated height,
\begin{equation}
\bigl[\mathbf{T}_{wc}\,(a\,\tilde{d}\,\mathbf{r}_{cam})\bigr]_z = Z,
\label{eq:had-constraint}
\end{equation}
which after multiplying out is linear in $(a, b)$. Stacking the
constraints for all instances in a camera into one over-determined
system gives a 2-unknown linear LSQ — solved in closed form, in
microseconds, without any iteration or learning. We refer the reader
to the supplementary for the explicit row form; the key practical
property is that the matrix conditioning is governed by the
\emph{diversity of object heights and depths} across the scene, not by
LiDAR density. A handful of cars + a few traffic-light poles is enough.


\subsection{Axis-Aligned HAD (AA-HAD)}
\label{sec:method-aahad}

The HAD constraints in Eq.~\ref{eq:had-constraint} involve the full
projection matrix $\mathbf{T}_{wc}$. For roadside cameras with
near-canonical mounting (tilt $\theta_t$ small, no roll), the same
constraints simplify to a per-camera \emph{axis-aligned} form in which
$z_{cam}$ becomes the world-Z primitive: the rotation
$\mathbf{T}_{wc}\!\!\downarrow_z$ acts on the camera-frame z-axis as
near identity, and the per-anchor constraint reduces to
$z_{cam}^{\text{anchor}} \approx \alpha\,(Z^{\text{anchor}} - t_z) + \beta$
for tilt-dependent $(\alpha, \beta)$. The full derivation is in the
supplementary. In code, AA-HAD is a single per-camera $\texttt{(a, b)}$
JSON file, written once and consumed by every downstream stage.

A second axis-alignment step happens at the \emph{pipeline} level
(Sec.~\ref{sec:method-pipeline}): after AA-HAD has put every pixel at
the right world Z, we snap road-surface pixels to a LiDAR-derived 2.5D
height grid. This snaps away the residual horizontal error that AA-HAD
cannot correct (vertical-axis-only correction by construction). The
combination — AA-HAD on the vertical axis + ground-snap on the
horizontal plane — is what we call the \emph{Layer~1+3 prior}, and it
is the input to both the LiDAR-priority fuse and the residual flow
head.


\subsection{The eight-stage pipeline}
\label{sec:method-pipeline}

Fig.~\ref{fig:pipeline} gives a graphical overview. The eight stages
correspond directly to the eight subcommands of the \texttt{lad} CLI
(Sec.~\ref{sec:impl}). Stages \texttt{depth}, \texttt{mask}, and
\texttt{seg} are off-the-shelf foundation-model runs — Depth Anything
3 \cite{depthanything3}, SAM \cite{sam}, SegFormer \cite{segformer} —
and produce per-frame artefacts that the geometry stages consume.
Stages \texttt{calib} and \texttt{object-accum} are the geometric
pre-processing; stages \texttt{complete}, \texttt{inject}, and
\texttt{render-bev} are the per-scene reconstruction and rendering.

\begin{figure*}[t]
  \centering
  \includegraphics[width=0.95\textwidth]{figures/pipeline.png}
  \caption{\textbf{The \ourmethod{} pipeline.} Eight stages, each with
  a typed config and its own per-run output directory. Solid arrows are
  data dependencies; dashed arrows are auto-discovery via the
  \texttt{outputs/<scene>/<upstream>/latest/} symlink convention.
  Multi-GPU sharding (\texttt{--runtime.gpu-ids}) is supported on every
  stage that benefits from it.}
  \label{fig:pipeline}
\end{figure*}


\subsection{Layer 3: ground-snap}
\label{sec:method-ground}

After AA-HAD has placed every pixel at a metric world position, we
build a 2.5D ground height grid $g(x, y)$ from the per-frame static
LiDAR returns (LiDAR points outside any V2X bbox), aggregated across
the scene. Each $(x, y)$ cell of the grid stores the low-quantile $Z$
of the LiDAR points that fell into it — the low quantile is robust
against occasional above-ground returns from grazing-incidence
reflectors (curbs, manhole covers, traffic-cone bases).

For each AA-HAD pixel with world coordinates $(x, y, z)$ and
$|z - g(x,y)| \le \delta_g$ (default $\delta_g = 0.4$\,m), we
\emph{snap} the pixel to $g(x, y)$. The residual flow head
(Sec.~\ref{sec:method-residual}) is responsible for any further
vertical adjustment within the ground band. The snap protects against
two well-known failure modes of AA-HAD: (i) per-camera tilt
miscalibration of a few tenths of a degree, which compounds over
50–100\,m of range, and (ii) DA3-introduced curvature near image
borders where the relative depth is least constrained.


\subsection{Layer 4: LiDAR-priority voxel fusion}
\label{sec:method-fuse}

Where the LiDAR is present, it should win: it is cm-precise, while
even a well-calibrated AA-HAD prior has ~5\,cm error. We voxelise both
clouds at the same grid (default 5\,cm), and for each voxel: if any
LiDAR point is present, keep \emph{only} the LiDAR points; otherwise,
keep the AA-HAD points (subject to an optional max-distance-to-nearest-
LiDAR cull that filters far-from-LiDAR AA-HAD outliers). The result
is what we call the \emph{hybrid cloud}: dense AA-HAD fills the
visible cones; LiDAR backbones the columns and the near range where
its rings are dense.

One operational nuance is worth calling out. If we keep the LiDAR
ground-band returns in the backbone, the BEV rendering inherits the
LiDAR sensor's \emph{ring pattern} — four concentric rings, one per
sensor, sweeping out from each sensor — because every ring voxel is
filled while the gaps between rings are sparser. Dropping the LiDAR
ground band (Sec.~\ref{sec:impl-skip-ground}) before the fuse lets the
uniform AA-HAD fill take over the road surface, giving a visually
uniform BEV.


\subsection{Layer 5: per-object accumulation}
\label{sec:method-perobject}

For each dynamic V2X object id, we accumulate its LiDAR + AA-HAD
returns across every frame in which the object appears, transforming
each contribution into the \emph{object-local} frame defined by that
frame's bounding box. The accumulated cloud is then voxel-downsampled
in the object frame and stored as a single PLY per object id. This
"freezes" the object's geometry in a canonical pose, so the per-ts
snapshot inject (Sec.~\ref{sec:method-inject}) only needs to look up
the object's pose at the current timestamp.

Per-object accumulation is also where the dataset's two V2X label
sources interact. The 10\,Hz interpolated boxes give us a thick stack
of per-frame observations per object — important because each
individual frame's bbox is too noisy to define a stable local frame on
its own. The 1\,Hz hand-labeled boxes are reserved for the
\emph{static-scene} accumulation (Layer 4), where bounding-box
precision is the limiting factor (Sec.~\ref{sec:impl-dual}).


\subsection{Residual flow head}
\label{sec:method-residual}

The Layer~1+3 prior is metric and roughly correct but exhibits
residual structure: edges where DA3's depth is uncertain, far-range
points where AA-HAD has a few cm of bias, and corners where the
ground-snap discontinuity meets vertical surfaces. We train a
point-cloud residual head $\mathbf{f}_\theta$ that maps the prior
cloud $\mathbf{P}\in\mathbb{R}^{N\times 3}$ to a per-point
displacement $\Delta\mathbf{P}\in\mathbb{R}^{N\times 3}$, and apply it
to obtain the refined cloud
$\mathbf{P}^* = \mathbf{P} + \mathbf{f}_\theta(\mathbf{P}, \mathbf{F})$.

\paragraph{Backbone.} We use a Dynamic Graph CNN \cite{dgcnn}
encoder with $k$-nearest-neighbour EdgeConv layers. Per-point features
$\mathbf{F}$ concatenate the RGB colour of the source pixel, a clipped
distance-to-nearest-LiDAR scalar, and a per-pixel Cityscapes class id
from SegFormer \cite{segformer} (the \emph{A1 feature}).

\paragraph{Segment-aware EdgeConv.} A naive EdgeConv KNN graph mixes
features across instance boundaries — e.g. a vehicle's points and the
road points next to it become neighbours. We mask the KNN graph to
exclude cross-segment edges using the Cityscapes class id, so messages
only flow within the same semantic class. This is the \emph{A1
masking} mentioned in the ablation.

\paragraph{Conditional flow-matching head.} The head outputs
$\Delta\mathbf{P}$ via a conditional flow-matching (CFM)
\cite{conditionalflowmatching,rectifiedflow} integrator. Compared to
a single-shot regression head, CFM is well-behaved in the small-Δxyz
regime that our task demands: the predicted Δxyz is usually on the
order of a few cm, and CFM's iterative refinement (default 4 steps)
avoids the regression collapse that a one-shot MSE head exhibits when
the target Δxyz is dominated by small displacements.

\paragraph{Per-frame deterministic seeding.} Both the random
subsample of input points and the CFM initial noise are seeded by a
hash of $(\text{scene}, \text{cam}, \text{ts})$, so a multi-GPU run
produces bit-for-bit identical refined clouds to a single-GPU run.


\subsection{Dynamic snapshot injection}
\label{sec:method-inject}

At a chosen anchor timestamp, we transform each per-object PLY from its
object-local frame back to the world frame using that timestamp's V2X
bounding-box pose. The result is concatenated onto the hybrid static
cloud, voxel-downsampled, and written as a single \texttt{hybrid\_with\_dyn.ply}.


\subsection{BEV video rendering with pose interpolation}
\label{sec:method-bev}

For video output, the snapshot inject runs per timestamp. At each
timestamp $t$, the per-object pose is obtained by linearly
interpolating the two annotated bracketing timestamps in the V2X
stream; objects whose annotation window is too far from $t$ (default
$\pm 200$\,ms) are not rendered. A separate check rejects
interpolation across long V2X tracker-loss gaps (default $> 500$\,ms),
which would otherwise teleport objects across the scene as the
tracker re-acquires them. The rendered cloud is rasterised top-down
with a z-buffer (highest world-Z wins per cell), producing one PNG per
timestamp. ffmpeg assembles the PNGs into a video.


%======================================================================
\section{Implementation}
\label{sec:impl}

\paragraph{Typed multi-stage CLI.} The entire pipeline ships as a
single \texttt{lad} command with eight subcommands:
\texttt{depth}, \texttt{mask}, \texttt{seg}, \texttt{calib},
\texttt{object-accum}, \texttt{complete}, \texttt{inject},
\texttt{render-bev}, plus a \texttt{pipeline} orchestrator that runs
all eight back-to-back. Configurations are typed \texttt{@dataclass}
trees consumed by \texttt{tyro}; per-stage artefacts land in
\texttt{outputs/<scene>/<stage>/<run\_id>/} with a sibling
\texttt{latest} symlink. Each downstream stage's input paths default
to \texttt{None}; the stage's \texttt{run()} call walks the
\texttt{latest} symlink of its declared upstream stage to fill them
in (Sec.~\ref{sec:impl-discovery}). The result is a one-shot
end-to-end invocation:

\begin{small}
\begin{verbatim}
lad pipeline \
  --complete.scene.scene 002 \
  --complete.scene.data-root /mnt/.../TianJin \
  --complete.runtime.gpu-ids 0 1 2 3 \
  --complete.residual-checkpoint ckpt.pt \
  --stages depth mask seg calib \
           object-accum complete \
           inject render-bev
\end{verbatim}
\end{small}


\subsection{Multi-GPU sharding}
\label{sec:impl-multigpu}

The per-frame stages (\texttt{depth}, \texttt{seg}) and the per-camera
stage (\texttt{mask}) shard their per-frame work across GPUs via the
\texttt{--runtime.gpu-ids} list. Each worker pins itself to one GPU
via \texttt{CUDA\_VISIBLE\_DEVICES} \emph{before} importing torch, so
that \texttt{cuda:0} resolves to the assigned card in the child
process; the parent never touches CUDA, avoiding the well-known
\texttt{spawn}-vs-\texttt{fork} deadlock that plagues
torch + multiprocessing pipelines. The \texttt{complete} stage shards
\texttt{(cam, ts)} work items via \texttt{torch.multiprocessing.Pool}.
The \texttt{object-accum} stage shards over the per-object loop via a
plain CPU \texttt{ProcessPoolExecutor} (16 workers by default).


\subsection{Auto-discovery of upstream artefacts}
\label{sec:impl-discovery}

Each downstream-stage config exposes its upstream-path fields as
\texttt{None} by default. At \texttt{run()} time, the stage calls a
small \texttt{engine.discovery} helper that walks
\texttt{<output.root>/<scene>/<upstream\_stage>/latest/} and fills in
the missing paths. Explicit overrides are still supported (a user
can pass \texttt{--complete.sam-mask-dir preview/sam/} to point at a
legacy product). The error path is also explicit: if the upstream
stage has never run, the discovery helper raises with a clear
migration hint (\texttt{run `lad <upstream>` first, or pass
--<flag> explicitly}).


\subsection{Dual labels source}
\label{sec:impl-dual}

The TianJin Roadside V2X dataset ships two parallel V2X annotation
sets per scene: \texttt{road\_labels/interpolation\_labels/} contains
the 10\,Hz interpolated bounding boxes; \texttt{road\_labels/
merged\_pcd\_all/} contains the 1\,Hz hand-labelled set from which the
interpolation is derived. The interpolation is convenient for
per-frame work (smooth across video) but its jitter at bbox edges
leaks dynamic-object LiDAR into the "static" accumulation: in a
moving vehicle's wake, the per-frame bbox is off by a few centimetres,
the vehicle's LiDAR returns are not fully culled, and across $N$
frames they form a faint trail (Fig.~\ref{fig:ablation-labels}). The
hand-labelled set is too sparse for video but its bbox precision is
much higher.

\ourmethod{} adopts a hybrid: \texttt{SceneConfig.static\_labels\_source
= "merged\_pcd"} drives the static-scene accumulation (high bbox
precision), and \texttt{SceneConfig.dynamic\_labels\_source =
"interpolation"} drives per-frame work (smooth motion). Both can be
overridden on the CLI if a particular scene lacks hand-labels.


\subsection{LiDAR ground-band skip}
\label{sec:impl-skip-ground}

The Layer 4 LiDAR-priority fuse normally keeps every in-camera-FOV
LiDAR point. As discussed in Sec.~\ref{sec:method-fuse}, this leaves
visible scan-ring artefacts on the road surface in the BEV. The
\texttt{--lidar-skip-ground} flag (on by default) drops LiDAR points
whose world Z is within $\pm 0.5$\,m of the LiDAR-derived ground
grid, \emph{before} they enter the fuse — the road geometry is already
locked by the ground-snap, so the cm-precision LiDAR ground points are
decorative. A subtlety: the outlier-rejection KD-tree inside
\texttt{lidar\_priority\_fill} must be built on the un-thinned LiDAR
backbone (otherwise the nearest-LiDAR query for AA-HAD ground points
returns several-metres-up vertical structures and the AA-HAD ground
gets falsely dropped, leaving a "halo around verticals" artefact in
the BEV).


%======================================================================
\section{Experiments}
\label{sec:exp}

\subsection{Setup}
\label{sec:exp-setup}

\paragraph{Dataset.} We use the TianJin Roadside V2X dataset, an
internal capture of 4 pinhole cameras (ids 0, 3, 6, 9) plus 4 LiDAR
sensors at a single intersection, with V2X 3D bounding boxes for every
visible vehicle, pedestrian, and roadside fixture. The release we
report on includes one scene used for training the residual head
(scene \texttt{008}, $\sim$190 frames per camera, $\sim$19\,s of
capture) and a second scene for zero-shot transfer (scene
\texttt{002}, similar size).

\paragraph{Models.} Depth Anything 3 (large variant)
\cite{depthanything3}, SAM (vit-huge) \cite{sam}, SegFormer (b5
fine-tuned on Cityscapes) \cite{segformer}. The residual flow head is
trained on scene \texttt{008} only, using the per-frame Layer~1+3
prior as input and the nearest static-LiDAR point as the regression
target.

\paragraph{Hardware.} 4$\times$NVIDIA A100 80\,GB, 100+ CPU cores,
1\,TB RAM. Total wall-clock time to take a fresh scene from raw V2X
data to a BEV video: roughly \todo{measure} minutes
(Sec.~\ref{sec:exp-runtime}).


\subsection{Closed-form vs.\ refined vs.\ hybrid}
\label{sec:exp-qualitative}

Fig.~\ref{fig:008-comparison} compares the three artefacts that the
\texttt{complete} stage produces on scene \texttt{008}:
\texttt{baseline.ply} (Layer 1+3 prior, no network),
\texttt{refined.ply} (baseline + residual flow head), and
\texttt{hybrid.ply} (refined + LiDAR-priority fuse). The closed-form
baseline alone already produces a recognisable scene; the residual
head cleans up edge bleed and removes residual ground curvature; the
LiDAR-priority fuse adds the cm-precise backbone needed for
metric-accurate downstream perception.


\begin{figure*}[t]
  \centering
  \includegraphics[width=0.95\textwidth]{figures/008_baseline_refined_hybrid.png}
  \caption{\textbf{Three artefacts of the \texttt{complete} stage on
  scene 008.} Left: closed-form AA-HAD baseline. Middle: refined by
  the residual flow head. Right: hybrid (refined + LiDAR-priority
  fuse). The refined cloud cleans up edge bleed and the residual
  ground curvature visible at the far end of the road in the baseline.
  The hybrid cloud adds the LiDAR backbone on poles, walls, and the
  vehicle silhouettes that LiDAR sees natively.}
  \label{fig:008-comparison}
\end{figure*}


\subsection{Zero-shot cross-scene transfer}
\label{sec:exp-zeroshot}

To probe whether the residual flow head generalises beyond its
training scene, we apply the \texttt{008}-trained checkpoint to scene
\texttt{002} \emph{without retraining}. The closed-form AA-HAD
calibration is re-fit per-camera on the new scene (Layer~1+3 is
inherently per-scene), but the network weights are frozen.
Fig.~\ref{fig:002-zeroshot} shows the result.

\begin{figure}[t]
  \centering
  \includegraphics[width=\linewidth]{figures/002_zeroshot.png}
  \caption{\textbf{Zero-shot transfer to an unseen scene.} The residual
  flow head trained on scene 008 is applied to scene 002. Top:
  closed-form baseline. Bottom: refined. The refined cloud is visibly
  cleaner than the baseline despite no exposure to 002 at training time
  — the residual head learned a generic geometric refinement, not a
  scene-specific one.}
  \label{fig:002-zeroshot}
\end{figure}


\subsection{Dynamic object reconstruction}
\label{sec:exp-dynamic}

Fig.~\ref{fig:per-object} shows three per-object accumulated PLYs
overlaid in the BEV at a chosen anchor timestamp. The per-object
clouds are denser than any single-frame snapshot would be — the LiDAR
contribution alone is too sparse for a usable vehicle silhouette — and
exhibit the front/rear/side asymmetry expected from a single
fly-through observation. The optional \texttt{--mirror-axis y} flag
(on for the result shown) doubles each object by reflecting across the
object-local $y$-axis, which fills in the occluded side at the cost of
a slight symmetry artefact.

\begin{figure}[t]
  \centering
  \includegraphics[width=\linewidth]{figures/per_object.png}
  \caption{\textbf{Per-object accumulated reconstructions placed in
  the BEV.} Three dynamic objects (two cars and a bus) at the anchor
  timestamp. The per-object accumulation runs in the object-local
  frame across every frame the object appears in, then transforms back
  to the world frame at the anchor.}
  \label{fig:per-object}
\end{figure}

A 4-frame video strip (Fig.~\ref{fig:bev-video}) demonstrates the
linear pose interpolation: a single vehicle is shown at $t$, $t{+}100$\,ms,
$t{+}200$\,ms, $t{+}300$\,ms. Between annotated V2X timestamps (which
fall on every 100\,ms boundary), the interpolated pose is geometrically
smooth and the vehicle's reconstruction translates correctly across
frames. A long V2X tracker-loss gap (Sec.~\ref{sec:method-bev}) would
manifest as the vehicle disappearing and reappearing at a distant
position; our $> 500$\,ms gap-cull suppresses this teleport.

\begin{figure*}[t]
  \centering
  \includegraphics[width=0.95\textwidth]{figures/bev_video_strip.png}
  \caption{\textbf{BEV video strip.} Four consecutive frames
  ($t = $0, 100, 200, 300\,ms) showing the same dynamic object placed
  in the scene via linear pose interpolation between annotated V2X
  timestamps. The reconstruction is anchored to the scene's static
  background (which is constant across frames) so only the dynamic
  layer moves.}
  \label{fig:bev-video}
\end{figure*}


\subsection{Ablations}
\label{sec:exp-ablation}

\paragraph{Ground-band LiDAR skip.}
Fig.~\ref{fig:ablation-skipground} shows a BEV close-up of the
intersection with and without the ground-band LiDAR skip. Without
the skip, the four LiDAR sensors' concentric scan rings are clearly
visible on the road surface; with the skip, the road density is
uniform.

\begin{figure}[t]
  \centering
  \includegraphics[width=\linewidth]{figures/ablation_skipground.png}
  \caption{\textbf{Ablation: LiDAR ground-band skip.} Left: every
  in-FOV LiDAR point participates in Layer~4 fuse. Right:
  \texttt{--lidar-skip-ground}: LiDAR points in the ground band are
  dropped before the fuse. The right panel is visibly more uniform.}
  \label{fig:ablation-skipground}
\end{figure}

\paragraph{Dual labels source.}
Fig.~\ref{fig:ablation-labels} compares the static cloud built from
interpolated 10\,Hz boxes vs.\ hand-labelled 1\,Hz boxes. The
interpolated build (left) shows clear "trails" along moving-vehicle
trajectories — interpolation jitter at the bbox edges lets a few
percent of each vehicle's LiDAR returns escape into the static branch
on every frame, and across $N$ frames the trails are visible. The
hand-labelled build (right) eliminates the trails at the cost of using
only the 1\,Hz subset of frames for accumulation.

\begin{figure}[t]
  \centering
  \includegraphics[width=\linewidth]{figures/ablation_labels.png}
  \caption{\textbf{Ablation: dual labels source.} Left: static cloud
  built from the 10\,Hz interpolated V2X boxes. Right: the same cloud
  built from the 1\,Hz hand-labelled boxes. Note the disappearance of
  the bus-trail and car-trail artefacts on the road surface (insets).
  The hand-labelled boxes are too sparse for per-frame dynamic-object
  rendering, so the released system uses each label source where it is
  strongest.}
  \label{fig:ablation-labels}
\end{figure}

\paragraph{Outlier reference for max-dist-to-LiDAR cull.}
A subtle interaction between the ground-band skip and the
\texttt{max\_dist\_to\_lidar} outlier filter inside
\texttt{lidar\_priority\_fill} is shown in
Fig.~\ref{fig:ablation-outlier-ref}. When the KD-tree underlying the
filter is built on the post-skip LiDAR, the AA-HAD road points'
nearest-LiDAR query returns a several-metres-up vertical structure;
those AA-HAD points are then falsely dropped, leaving a "halo around
verticals" pattern. Building the KD-tree on the un-thinned LiDAR
(while the backbone still uses the post-skip subset) fixes the
artefact.

\begin{figure}[t]
  \centering
  \includegraphics[width=\linewidth]{figures/ablation_outlier_ref.png}
  \caption{\textbf{Ablation: outlier-rejection reference cloud.} The
  \texttt{max-dist-to-lidar} cull keeps AA-HAD points that are close
  to a LiDAR return. When \texttt{lidar\_priority\_fill}'s KD-tree is
  built on the post-skip LiDAR (left), the cull falsely drops AA-HAD
  road points (the nearest LiDAR is now a vertical structure several
  metres up). Pointing the KD-tree at the un-thinned LiDAR while still
  using the post-skip subset for the backbone (right) gives a clean
  road surface.}
  \label{fig:ablation-outlier-ref}
\end{figure}


\subsection{Runtime}
\label{sec:exp-runtime}

Table~\ref{tab:runtime} reports per-stage wall-clock times on the
4$\times$A100 + 100-core hardware described in
Sec.~\ref{sec:exp-setup}. The hot path is per-(cam, ts) network
refinement in the \texttt{complete} stage, parallelised across the 4
GPUs via \texttt{torch.multiprocessing.Pool}. The
\texttt{object-accum} stage runs CPU-only with 16 parallel
\texttt{ProcessPoolExecutor} workers.

\begin{table}[t]
  \centering
  \small
  \begin{tabular}{lr@{\,}lr}
    \toprule
    Stage & \multicolumn{2}{c}{Hardware} & Wall time \\
    \midrule
    \texttt{depth}        & 4$\times$ & A100 & $\sim$3\,min \\
    \texttt{mask}         & 4$\times$ & A100 & $\sim$2\,min \\
    \texttt{seg}          & 4$\times$ & A100 & $\sim$2\,min \\
    \texttt{calib}        & 1$\times$ & A100 & $\sim$2\,min \\
    \texttt{object-accum} & 16$\times$ & CPU & $\sim$30--60\,min \\
    \texttt{complete}     & 4$\times$ & A100 & $\sim$30\,min \\
    \texttt{inject}       & 1$\times$ & CPU & $\sim$2\,min \\
    \texttt{render-bev}   & 1$\times$ & CPU & $\sim$5--10\,min \\
    \midrule
    \textbf{Total} (\texttt{pipeline}) & & & $\sim$75--120\,min \\
    \bottomrule
  \end{tabular}
  \caption{\textbf{End-to-end runtime for one scene} ($\sim$190 frames per
  camera, 4 cameras, $\sim$20 dynamic objects). All stages are
  resumable: \texttt{--from-stage <X>} skips upstream stages that
  succeeded, \texttt{--skip-existing} (on by default in
  \texttt{object-accum}) lets a killed run pick up where it left off.}
  \label{tab:runtime}
\end{table}


%======================================================================
\section{Discussion and limitations}
\label{sec:discussion}

\paragraph{What HAD assumes.} HAD assumes (i) per-instance bounding
boxes are available — a reasonable assumption in V2X but not in the
general monocular setting; (ii) the road plane is locally
approximately level — which holds for urban intersections but degrades
on steep grades; and (iii) the camera tilt is well-calibrated to a
fraction of a degree — which is the AA-HAD axis-alignment assumption,
and which the ground-snap covers for small tilt errors.

\paragraph{Residual head generalisation.} Our zero-shot transfer
result (Sec.~\ref{sec:exp-zeroshot}) shows the residual head trained
on one intersection generalises to another, but this is a single
data point. We expect the head to need re-training (or fine-tuning)
when applied to scenes with substantially different camera mounting
heights, sensor configurations, or scene types (e.g.\ rural roads,
parking lots).

\paragraph{Dynamic-object completeness.} The per-object accumulated
clouds inherit the self-occlusion of the single fly-through
observation: a vehicle passing on the far side of the intersection is
observed only from one side, and the reconstructed object is
missing the occluded faces. The \texttt{--mirror-axis y} flag
mitigates this for vehicles assumed to be left-right symmetric, but a
principled fix would use a learnt object prior — the natural next
extension of the system.


%======================================================================
\section{Conclusion}
\label{sec:conclusion}

We introduced \ourmethod{}, a closed-form-anchored dense 3D
reconstruction pipeline for roadside V2X scenes. The conceptual core
is the Height-Anchored Depth formulation, which calibrates a
foundation depth model against per-instance world-frame heights in a
single one-shot linear solve. Around this core we have built a
typed eight-stage CLI that turns raw V2X data into both dense static
reconstructions and BEV videos with smooth dynamic actors, in $\sim$2
hours of wall time per scene on four A100s.

The two design choices we hope generalise beyond this specific
pipeline are (i) \emph{anchor on what is annotated} — per-instance
heights, not per-pixel depths; the V2X data already gives us the
former — and (ii) \emph{anchor in the world frame}, where the scene
geometry is well-conditioned, rather than in the camera frame where
foundation models naturally produce relative outputs.


%======================================================================
{\small
\bibliographystyle{ieee_fullname}
\bibliography{references}
}

\end{document}
